% Options for packages loaded elsewhere
\PassOptionsToPackage{unicode}{hyperref}
\PassOptionsToPackage{hyphens}{url}
\PassOptionsToPackage{dvipsnames,svgnames,x11names}{xcolor}
%
\documentclass[
  super,
  preprint,
  3p]{elsarticle}

\usepackage{amsmath,amssymb}
\usepackage{iftex}
\ifPDFTeX
  \usepackage[T1]{fontenc}
  \usepackage[utf8]{inputenc}
  \usepackage{textcomp} % provide euro and other symbols
\else % if luatex or xetex
  \usepackage{unicode-math}
  \defaultfontfeatures{Scale=MatchLowercase}
  \defaultfontfeatures[\rmfamily]{Ligatures=TeX,Scale=1}
\fi
\usepackage{lmodern}
\ifPDFTeX\else  
    % xetex/luatex font selection
\fi
% Use upquote if available, for straight quotes in verbatim environments
\IfFileExists{upquote.sty}{\usepackage{upquote}}{}
\IfFileExists{microtype.sty}{% use microtype if available
  \usepackage[]{microtype}
  \UseMicrotypeSet[protrusion]{basicmath} % disable protrusion for tt fonts
}{}
\makeatletter
\@ifundefined{KOMAClassName}{% if non-KOMA class
  \IfFileExists{parskip.sty}{%
    \usepackage{parskip}
  }{% else
    \setlength{\parindent}{0pt}
    \setlength{\parskip}{6pt plus 2pt minus 1pt}}
}{% if KOMA class
  \KOMAoptions{parskip=half}}
\makeatother
\usepackage{xcolor}
\setlength{\emergencystretch}{3em} % prevent overfull lines
\setcounter{secnumdepth}{5}
% Make \paragraph and \subparagraph free-standing
\makeatletter
\ifx\paragraph\undefined\else
  \let\oldparagraph\paragraph
  \renewcommand{\paragraph}{
    \@ifstar
      \xxxParagraphStar
      \xxxParagraphNoStar
  }
  \newcommand{\xxxParagraphStar}[1]{\oldparagraph*{#1}\mbox{}}
  \newcommand{\xxxParagraphNoStar}[1]{\oldparagraph{#1}\mbox{}}
\fi
\ifx\subparagraph\undefined\else
  \let\oldsubparagraph\subparagraph
  \renewcommand{\subparagraph}{
    \@ifstar
      \xxxSubParagraphStar
      \xxxSubParagraphNoStar
  }
  \newcommand{\xxxSubParagraphStar}[1]{\oldsubparagraph*{#1}\mbox{}}
  \newcommand{\xxxSubParagraphNoStar}[1]{\oldsubparagraph{#1}\mbox{}}
\fi
\makeatother

\usepackage{color}
\usepackage{fancyvrb}
\newcommand{\VerbBar}{|}
\newcommand{\VERB}{\Verb[commandchars=\\\{\}]}
\DefineVerbatimEnvironment{Highlighting}{Verbatim}{commandchars=\\\{\}}
% Add ',fontsize=\small' for more characters per line
\usepackage{framed}
\definecolor{shadecolor}{RGB}{241,243,245}
\newenvironment{Shaded}{\begin{snugshade}}{\end{snugshade}}
\newcommand{\AlertTok}[1]{\textcolor[rgb]{0.68,0.00,0.00}{#1}}
\newcommand{\AnnotationTok}[1]{\textcolor[rgb]{0.37,0.37,0.37}{#1}}
\newcommand{\AttributeTok}[1]{\textcolor[rgb]{0.40,0.45,0.13}{#1}}
\newcommand{\BaseNTok}[1]{\textcolor[rgb]{0.68,0.00,0.00}{#1}}
\newcommand{\BuiltInTok}[1]{\textcolor[rgb]{0.00,0.23,0.31}{#1}}
\newcommand{\CharTok}[1]{\textcolor[rgb]{0.13,0.47,0.30}{#1}}
\newcommand{\CommentTok}[1]{\textcolor[rgb]{0.37,0.37,0.37}{#1}}
\newcommand{\CommentVarTok}[1]{\textcolor[rgb]{0.37,0.37,0.37}{\textit{#1}}}
\newcommand{\ConstantTok}[1]{\textcolor[rgb]{0.56,0.35,0.01}{#1}}
\newcommand{\ControlFlowTok}[1]{\textcolor[rgb]{0.00,0.23,0.31}{\textbf{#1}}}
\newcommand{\DataTypeTok}[1]{\textcolor[rgb]{0.68,0.00,0.00}{#1}}
\newcommand{\DecValTok}[1]{\textcolor[rgb]{0.68,0.00,0.00}{#1}}
\newcommand{\DocumentationTok}[1]{\textcolor[rgb]{0.37,0.37,0.37}{\textit{#1}}}
\newcommand{\ErrorTok}[1]{\textcolor[rgb]{0.68,0.00,0.00}{#1}}
\newcommand{\ExtensionTok}[1]{\textcolor[rgb]{0.00,0.23,0.31}{#1}}
\newcommand{\FloatTok}[1]{\textcolor[rgb]{0.68,0.00,0.00}{#1}}
\newcommand{\FunctionTok}[1]{\textcolor[rgb]{0.28,0.35,0.67}{#1}}
\newcommand{\ImportTok}[1]{\textcolor[rgb]{0.00,0.46,0.62}{#1}}
\newcommand{\InformationTok}[1]{\textcolor[rgb]{0.37,0.37,0.37}{#1}}
\newcommand{\KeywordTok}[1]{\textcolor[rgb]{0.00,0.23,0.31}{\textbf{#1}}}
\newcommand{\NormalTok}[1]{\textcolor[rgb]{0.00,0.23,0.31}{#1}}
\newcommand{\OperatorTok}[1]{\textcolor[rgb]{0.37,0.37,0.37}{#1}}
\newcommand{\OtherTok}[1]{\textcolor[rgb]{0.00,0.23,0.31}{#1}}
\newcommand{\PreprocessorTok}[1]{\textcolor[rgb]{0.68,0.00,0.00}{#1}}
\newcommand{\RegionMarkerTok}[1]{\textcolor[rgb]{0.00,0.23,0.31}{#1}}
\newcommand{\SpecialCharTok}[1]{\textcolor[rgb]{0.37,0.37,0.37}{#1}}
\newcommand{\SpecialStringTok}[1]{\textcolor[rgb]{0.13,0.47,0.30}{#1}}
\newcommand{\StringTok}[1]{\textcolor[rgb]{0.13,0.47,0.30}{#1}}
\newcommand{\VariableTok}[1]{\textcolor[rgb]{0.07,0.07,0.07}{#1}}
\newcommand{\VerbatimStringTok}[1]{\textcolor[rgb]{0.13,0.47,0.30}{#1}}
\newcommand{\WarningTok}[1]{\textcolor[rgb]{0.37,0.37,0.37}{\textit{#1}}}

\providecommand{\tightlist}{%
  \setlength{\itemsep}{0pt}\setlength{\parskip}{0pt}}\usepackage{longtable,booktabs,array}
\usepackage{calc} % for calculating minipage widths
% Correct order of tables after \paragraph or \subparagraph
\usepackage{etoolbox}
\makeatletter
\patchcmd\longtable{\par}{\if@noskipsec\mbox{}\fi\par}{}{}
\makeatother
% Allow footnotes in longtable head/foot
\IfFileExists{footnotehyper.sty}{\usepackage{footnotehyper}}{\usepackage{footnote}}
\makesavenoteenv{longtable}
\usepackage{graphicx}
\makeatletter
\newsavebox\pandoc@box
\newcommand*\pandocbounded[1]{% scales image to fit in text height/width
  \sbox\pandoc@box{#1}%
  \Gscale@div\@tempa{\textheight}{\dimexpr\ht\pandoc@box+\dp\pandoc@box\relax}%
  \Gscale@div\@tempb{\linewidth}{\wd\pandoc@box}%
  \ifdim\@tempb\p@<\@tempa\p@\let\@tempa\@tempb\fi% select the smaller of both
  \ifdim\@tempa\p@<\p@\scalebox{\@tempa}{\usebox\pandoc@box}%
  \else\usebox{\pandoc@box}%
  \fi%
}
% Set default figure placement to htbp
\def\fps@figure{htbp}
\makeatother

\makeatletter
\@ifpackageloaded{caption}{}{\usepackage{caption}}
\AtBeginDocument{%
\ifdefined\contentsname
  \renewcommand*\contentsname{Table of contents}
\else
  \newcommand\contentsname{Table of contents}
\fi
\ifdefined\listfigurename
  \renewcommand*\listfigurename{List of Figures}
\else
  \newcommand\listfigurename{List of Figures}
\fi
\ifdefined\listtablename
  \renewcommand*\listtablename{List of Tables}
\else
  \newcommand\listtablename{List of Tables}
\fi
\ifdefined\figurename
  \renewcommand*\figurename{Figure}
\else
  \newcommand\figurename{Figure}
\fi
\ifdefined\tablename
  \renewcommand*\tablename{Table}
\else
  \newcommand\tablename{Table}
\fi
}
\@ifpackageloaded{float}{}{\usepackage{float}}
\floatstyle{ruled}
\@ifundefined{c@chapter}{\newfloat{codelisting}{h}{lop}}{\newfloat{codelisting}{h}{lop}[chapter]}
\floatname{codelisting}{Listing}
\newcommand*\listoflistings{\listof{codelisting}{List of Listings}}
\makeatother
\makeatletter
\makeatother
\makeatletter
\@ifpackageloaded{caption}{}{\usepackage{caption}}
\@ifpackageloaded{subcaption}{}{\usepackage{subcaption}}
\makeatother
\journal{Journal Name}

\usepackage[]{natbib}
\bibliographystyle{elsarticle-num}
\usepackage{bookmark}

\IfFileExists{xurl.sty}{\usepackage{xurl}}{} % add URL line breaks if available
\urlstyle{same} % disable monospaced font for URLs
\hypersetup{
  pdftitle={Regresión logística un ejemplo para la predicción de infartos},
  pdfauthor={Alice Anonymous; Bob Security; Cat Memes; Derek Zoolander},
  pdfkeywords={keyword1, keyword2},
  colorlinks=true,
  linkcolor={blue},
  filecolor={Maroon},
  citecolor={Blue},
  urlcolor={Blue},
  pdfcreator={LaTeX via pandoc}}


\setlength{\parindent}{6pt}
\begin{document}

\begin{frontmatter}
\title{Regresión logística un ejemplo para la predicción de infartos}
\author[1]{Alice Anonymous%
\corref{cor1}%
\fnref{fn1}}
 \ead{alice@example.com} 
\author[2]{Bob Security%
%
\fnref{fn2}}
 \ead{bob@example.com} 
\author[2]{Cat Memes%
%
\fnref{fn3}}
 \ead{cat@example.com} 
\author[]{Derek Zoolander%
%
}
 \ead{derek@example.com} 

\affiliation[1]{organization={Some Institute of Technology, Department
Name},addressline={Street Address},city={City},postcode={Postal
Code},postcodesep={}}
\affiliation[2]{organization={Another University, Department
Name},addressline={Street Address},city={City},postcode={Postal
Code},postcodesep={}}

\cortext[cor1]{Corresponding author}
\fntext[fn1]{This is the first author footnote.}
\fntext[fn2]{Another author footnote, this is a very long footnote and
it should be a really long footnote. But this footnote is not yet
sufficiently long enough to make two lines of footnote text.}
\fntext[fn3]{Yet another author footnote.}

        
\begin{abstract}
This is the abstract. Lorem ipsum dolor sit amet, consectetur adipiscing
elit. Vestibulum augue turpis, dictum non malesuada a, volutpat eget
velit. Nam placerat turpis purus, eu tristique ex tincidunt et. Mauris
sed augue eget turpis ultrices tincidunt. Sed et mi in leo porta
egestas. Aliquam non laoreet velit. Nunc quis ex vitae eros aliquet
auctor nec ac libero. Duis laoreet sapien eu mi luctus, in bibendum leo
molestie. Sed hendrerit diam diam, ac dapibus nisl volutpat vitae.
Aliquam bibendum varius libero, eu efficitur justo rutrum at. Sed at
tempus elit.
\end{abstract}





\begin{keyword}
    keyword1 \sep 
    keyword2
\end{keyword}
\end{frontmatter}
    

\subsection{Capítulo 1
Introducción}\label{capuxedtulo-1-introducciuxf3n}

{[}cite\_start{]}Las \textbf{redes neuronales convolucionales profundas
(CNNs)} han estado en el centro de espectaculares avances en el
aprendizaje profundo{[}cite: 312{]}. {[}cite\_start{]}Aunque las CNNs se
han utilizado tan pronto como en los años noventa para resolver tareas
de reconocimiento de caracteres (Le Cun et al., 1997) {[}cite:
313{]}{[}cite\_start{]}, su aplicación generalizada actual se debe a
trabajos mucho más recientes, cuando se utilizó una CNN profunda para
superar el estado del arte en el desafío de clasificación de imágenes
ImageNet (Krizhevsky et al., 2012){[}cite: 313{]}.

{[}cite\_start{]}Por lo tanto, las redes neuronales convolucionales
constituyen una herramienta muy útil para los profesionales del
aprendizaje automático{[}cite: 314{]}. {[}cite\_start{]}Sin embargo,
aprender a usar las CNNs por primera vez es generalmente una experiencia
intimidante{[}cite: 315{]}. {[}cite\_start{]}La \textbf{forma de la
salida de una capa convolucional} se ve afectada por la forma de su
entrada, así como por la elección de la forma del núcleo
(\emph{kernel}), el relleno de ceros (\emph{zero padding}) y los pasos
(\emph{strides}), y la relación entre estas propiedades no es trivial de
inferir{[}cite: 316{]}. {[}cite\_start{]}Esto contrasta con las
\textbf{capas totalmente conectadas}, cuyo tamaño de salida es
independiente del tamaño de la entrada{[}cite: 317{]}.
{[}cite\_start{]}Además, las CNNs también suelen presentar una etapa de
\textbf{agrupamiento (\emph{pooling})}, lo que añade otro nivel de
complejidad con respecto a las redes totalmente conectadas{[}cite:
318{]}. {[}cite\_start{]}Finalmente, las llamadas \textbf{capas
convolucionales transpuestas} (también conocidas como capas
convolucionales con pasos fraccionados) se han empleado en cada vez más
trabajos últimamente (Zeiler et al., 2011; Zeiler y Fergus, 2014; Long
et al., 2015; Radford et al., 2015; Visin et al., 2015; Im et al., 2016)
{[}cite: 319{]}{[}cite\_start{]}, y su relación con las capas
convolucionales se ha explicado con diversos grados de claridad{[}cite:
319{]}.

{[}cite\_start{]}El objetivo de esta guía es doble{[}cite: 320{]}:

\begin{enumerate}
\def\labelenumi{\arabic{enumi}.}
\tightlist
\item
  {[}cite\_start{]}Explicar la relación entre las capas convolucionales
  y las capas convolucionales transpuestas{[}cite: 321{]}.
\item
  {[}cite\_start{]}Proporcionar una comprensión intuitiva de la relación
  entre la forma de la entrada, la forma del núcleo, el relleno de
  ceros, los pasos y la forma de la salida en las capas convolucionales,
  de \emph{pooling} y convolucionales transpuestas{[}cite: 322{]}.
\end{enumerate}

{[}cite\_start{]}Para seguir siendo ampliamente aplicables, los
resultados mostrados en esta guía son independientes de los detalles de
implementación y se aplican a todos los \emph{frameworks} de aprendizaje
automático de uso común, como Theano (Bergstra et al., 2010; Bastien et
al., 2012) {[}cite: 323{]}{[}cite\_start{]}, Torch (Collobert et al.,
2011), Tensorflow (Abadi et al., 2015) y Caffe (Jia et al.,
2014){[}cite: 325{]}.

{[}cite\_start{]}Este capítulo revisa brevemente los principales bloques
de construcción de las CNNs, a saber, las convoluciones discretas y el
\emph{pooling}{[}cite: 326{]}. {[}cite\_start{]}Para un tratamiento en
profundidad del tema, consulte el Capítulo 9 del libro de texto
\emph{Deep Learning} (Goodfellow et al., 2016){[}cite: 327{]}.

\begin{center}\rule{0.5\linewidth}{0.5pt}\end{center}

\subsubsection{1.1 Convoluciones
discretas}\label{convoluciones-discretas}

{[}cite\_start{]}El pan de cada día de las redes neuronales son las
\textbf{transformaciones afines}: se recibe un vector como entrada y se
multiplica por una matriz para producir una salida (a la que se suele
sumar un vector de sesgo antes de pasar el resultado a través de una no
linealidad){[}cite: 330{]}. {[}cite\_start{]}Esto es aplicable a
cualquier tipo de entrada, ya sea una imagen, un clip de sonido o una
colección desordenada de características: sea cual sea su
dimensionalidad, su representación siempre se puede aplanar en un vector
antes de la transformación{[}cite: 331{]}.

{[}cite\_start{]}Las imágenes, los clips de sonido y muchos otros tipos
de datos similares tienen una \textbf{estructura intrínseca}{[}cite:
332{]}. {[}cite\_start{]}Más formalmente, comparten estas propiedades
importantes{[}cite: 333{]}:

\begin{itemize}
\tightlist
\item
  {[}cite\_start{]}Se almacenan como \textbf{matrices
  multidimensionales}{[}cite: 334{]}.
\item
  {[}cite\_start{]}Presentan uno o más ejes para los que el
  \textbf{orden es importante} (por ejemplo, los ejes de ancho y alto
  para una imagen, el eje del tiempo para un clip de sonido){[}cite:
  335{]}.
\item
  {[}cite\_start{]}Se utiliza un eje, llamado \textbf{eje de canal},
  para acceder a diferentes vistas de los datos (por ejemplo, los
  canales rojo, verde y azul de una imagen en color, o los canales
  izquierdo y derecho de una pista de audio estéreo){[}cite: 336{]}.
\end{itemize}

{[}cite\_start{]}Estas propiedades no se explotan cuando se aplica una
transformación afín {[}cite: 337{]}{[}cite\_start{]}; de hecho, todos
los ejes se tratan de la misma manera y la información topológica no se
tiene en cuenta{[}cite: 338{]}. {[}cite\_start{]}Aún así, aprovechar la
estructura implícita de los datos puede resultar muy útil para resolver
algunas tareas, como la visión por computadora y el reconocimiento de
voz, y en estos casos sería mejor preservarla{[}cite: 339{]}.
{[}cite\_start{]}Aquí es donde entran en juego las \textbf{convoluciones
discretas}{[}cite: 340{]}.

{[}cite\_start{]}Una convolución discreta es una \textbf{transformación
lineal que preserva esta noción de orden}{[}cite: 341{]}.
{[}cite\_start{]}Es \textbf{dispersa} (solo unas pocas unidades de
entrada contribuyen a una unidad de salida dada) y \textbf{reutiliza
parámetros} (los mismos pesos se aplican a múltiples ubicaciones en la
entrada){[}cite: 342{]}.

{[}cite\_start{]}La Figura 1.1 proporciona un ejemplo de una convolución
discreta{[}cite: 343{]}. {[}cite\_start{]}La cuadrícula de color azul
claro se denomina \textbf{mapa de características de entrada}
(\emph{input feature map}){[}cite: 343{]}. Para simplificar el dibujo,
se representa un solo mapa de características de entrada, pero no es
raro tener múltiples mapas de características apilados uno encima del
otro. Un \textbf{núcleo} (\emph{kernel}) (área sombreada) de valor

\[
\begin{pmatrix}
2 & 2 & 3 \\
2 & 0 & 0 \\
2 & 0 & 2
\end{pmatrix}
\]

{[}cite\_start{]}se desliza a través del mapa de características de
entrada{[}cite: 356{]}. {[}cite\_start{]}En cada ubicación, se calcula
el producto entre cada elemento del núcleo y el elemento de entrada que
se superpone, y los resultados se suman para obtener la salida en la
ubicación actual{[}cite: 356{]}. {[}cite\_start{]}El procedimiento se
puede repetir utilizando diferentes núcleos para formar tantos mapas de
características de salida como se desee (Figura 1.3){[}cite: 357{]}.
{[}cite\_start{]}Las salidas finales de este procedimiento se denominan
\textbf{mapas de características de salida} (\emph{output feature
maps}){[}cite: 358{]}. {[}cite\_start{]}Si hay múltiples mapas de
características de entrada, el núcleo tendrá que ser tridimensional o,
de manera equivalente, cada uno de los mapas de características se
convolucionará con un núcleo distinto y los mapas de características
resultantes se sumarán elemento a elemento para producir el mapa de
características de salida{[}cite: 359{]}.

{[}cite\_start{]}La convolución representada en la Figura 1.1 es una
instancia de una \textbf{convolución 2-D}, pero se puede generalizar a
\textbf{convoluciones N-D}{[}cite: 360{]}. {[}cite\_start{]}Por ejemplo,
en una convolución 3-D, el núcleo sería un cuboide y se deslizaría a
través de la altura, el ancho y la profundidad del mapa de
características de entrada{[}cite: 361{]}.

{[}cite\_start{]}La colección de núcleos que definen una convolución
discreta tiene una forma que corresponde a alguna permutación de
\[(n, m, k_1, \dots, k_N)\] donde{[}cite: 362{]}:

\begin{itemize}
\tightlist
\item
  {[}cite\_start{]}\(n \equiv\) número de mapas de características de
  salida{[}cite: 363{]},
\item
  {[}cite\_start{]}\(m \equiv\) número de mapas de características de
  entrada{[}cite: 364{]},
\item
  {[}cite\_start{]}\(k_j \equiv\) tamaño del núcleo a lo largo del eje
  \(j\){[}cite: 364{]}.
\end{itemize}

{[}cite\_start{]}Las siguientes propiedades afectan el tamaño de salida
\(o_j\) de una capa convolucional a lo largo del eje \(j\){[}cite:
365{]}:

\begin{itemize}
\tightlist
\item
  {[}cite\_start{]}\(i_j\): tamaño de entrada a lo largo del eje
  \(j\){[}cite: 366{]},
\item
  {[}cite\_start{]}\(k_j\): tamaño del núcleo a lo largo del eje
  \(j\){[}cite: 367{]},
\item
  {[}cite\_start{]}\(s_j\): paso (distancia entre dos posiciones
  consecutivas del núcleo) a lo largo del eje \(j\){[}cite: 368{]},
\item
  {[}cite\_start{]}\(p_j\): relleno de ceros (número de ceros
  concatenados al principio y al final de un eje) a lo largo del eje
  \(j\){[}cite: 369{]}.
\end{itemize}

{[}cite\_start{]}Por ejemplo, la Figura 1.2 muestra un núcleo de
\(3 \times 3\) aplicado a una entrada de \(5 \times 5\) rellenada con un
borde de ceros de \(1 \times 1\) utilizando pasos de
\(2 \times 2\){[}cite: 370{]}.

{[}cite\_start{]}Tenga en cuenta que los pasos constituyen una forma de
\textbf{submuestreo}{[}cite: 371{]}. {[}cite\_start{]}Como alternativa a
ser interpretados como una medida de cuánto se traslada el núcleo, los
pasos también pueden verse como la cantidad de salida que se
retiene{[}cite: 371{]}. {[}cite\_start{]}Por ejemplo, mover el núcleo en
saltos de dos es equivalente a mover el núcleo en saltos de uno, pero
reteniendo solo los elementos de salida impares (Figura 1.4){[}cite:
372{]}.

{[}cite\_start{]}\textsuperscript{1}Un ejemplo de esto es lo que se
denominó anteriormente canales para imágenes y clips de sonido{[}cite:
373{]}.

{[}cite\_start{]}\textsuperscript{2}Aunque hay una distinción entre
convolución y correlación cruzada desde una perspectiva de procesamiento
de señales, las dos se vuelven intercambiables cuando se aprende el
núcleo{[}cite: 374{]}. {[}cite\_start{]}Por el bien de la simplicidad y
para ser coherentes con la mayor parte de la literatura de aprendizaje
automático, el término \textbf{convolución} se utilizará en esta
guía{[}cite: 375{]}.

\begin{center}\rule{0.5\linewidth}{0.5pt}\end{center}

\textbf{Figura 1.1:} Cálculo de los valores de salida de una convolución
discreta.

\textbf{Figura 1.2:} Cálculo de los valores de salida de una convolución
discreta para \(N=2\), \(i_1=i_2=5\), \(k_1=k_2=3\), \(s_1=s_2=2\), y
\(p_1=p_2=1\).

{[}cite\_start{]}\textbf{Figura 1.3:} Una convolución que mapea de dos
mapas de características de entrada a tres mapas de características de
salida utilizando una colección de núcleos \textbf{w} de
\(3 \times 2 \times 3 \times 3\){[}cite: 377{]}. {[}cite\_start{]}En la
trayectoria izquierda, el mapa de características de entrada 1 se
convoluciona con el núcleo \(w_{1,1}\) y el mapa de características de
entrada 2 se convoluciona con el núcleo \(w_{1,2}\), y los resultados se
suman elemento a elemento para formar el primer mapa de características
de salida{[}cite: 378{]}. {[}cite\_start{]}Se repite lo mismo para las
trayectorias central y derecha para formar el segundo y tercer mapa de
características, y los tres mapas de características de salida se
agrupan para formar la salida{[}cite: 379{]}.

\textbf{Figura 1.4:} Una forma alternativa de ver los pasos.
{[}cite\_start{]}En lugar de trasladar el núcleo de \(3 \times 3\) en
incrementos de \(s=2\) (izquierda), el núcleo se traslada en incrementos
de 1 y solo se retiene uno de cada \(s=2\) elementos de salida
(derecha){[}cite: 380{]}.

\begin{center}\rule{0.5\linewidth}{0.5pt}\end{center}

\subsubsection{1.2 Pooling (Agrupamiento)}\label{pooling-agrupamiento}

{[}cite\_start{]}Además de las convoluciones discretas en sí, las
\textbf{operaciones de \emph{pooling}} (agrupamiento) constituyen otro
bloque de construcción importante en las CNNs{[}cite: 383{]}.
{[}cite\_start{]}Las operaciones de \emph{pooling} reducen el tamaño de
los mapas de características utilizando alguna función para resumir
subregiones, como tomar el promedio o el valor máximo{[}cite: 384{]}.

{[}cite\_start{]}El \emph{Pooling} funciona deslizando una ventana a
través de la entrada y alimentando el contenido de la ventana a una
función de \emph{pooling}{[}cite: 385{]}. {[}cite\_start{]}En cierto
sentido, el \emph{pooling} funciona de manera muy similar a una
convolución discreta, pero reemplaza la combinación lineal descrita por
el núcleo con alguna otra función{[}cite: 386{]}. {[}cite\_start{]}La
Figura 1.5 proporciona un ejemplo de \emph{average pooling}
(agrupamiento promedio), y la Figura 1.6 hace lo mismo para \emph{max
pooling} (agrupamiento máximo){[}cite: 387{]}.

{[}cite\_start{]}Las siguientes propiedades afectan el tamaño de salida
\(o_j\) de una capa de \emph{pooling} a lo largo del eje \(j\){[}cite:
389{]}:

\begin{itemize}
\tightlist
\item
  {[}cite\_start{]}\(i_j\): tamaño de entrada a lo largo del eje
  \(j\){[}cite: 390{]},
\item
  {[}cite\_start{]}\(k_j\): tamaño de la ventana de \emph{pooling} a lo
  largo del eje \(j\){[}cite: 392{]},
\item
  {[}cite\_start{]}\(s_j\): paso (distancia entre dos posiciones
  consecutivas de la ventana de \emph{pooling}) a lo largo del eje
  \(j\){[}cite: 393{]}.
\end{itemize}

\begin{center}\rule{0.5\linewidth}{0.5pt}\end{center}

{[}cite\_start{]}\textbf{Figura 1.5:} Cálculo de los valores de salida
de una operación de \emph{average pooling} de \(3 \times 3\) en una
entrada de \(5 \times 5\) utilizando pasos de \(1 \times 1\){[}cite:
395{]}.

{[}cite\_start{]}\textbf{Figura 1.6:} Cálculo de los valores de salida
de una operación de \emph{max pooling} de \(3 \times 3\) en una entrada
de \(5 \times 5\) utilizando pasos de \(1 \times 1\){[}cite: 396{]}.

\begin{center}\rule{0.5\linewidth}{0.5pt}\end{center}

\subsection{Capítulo 2 Aritmética de la
convolución}\label{capuxedtulo-2-aritmuxe9tica-de-la-convoluciuxf3n}

{[}cite\_start{]}El análisis de la relación entre las propiedades de la
capa convolucional se facilita por el hecho de que \textbf{no
interactúan a través de los ejes} {[}cite: 398{]}{[}cite\_start{]}, es
decir, la elección del tamaño del núcleo, el paso y el relleno de ceros
a lo largo del eje \(j\) solo afecta el tamaño de la salida del eje
\(j\){[}cite: 398{]}.

{[}cite\_start{]}Debido a eso, este capítulo se centrará en el siguiente
ajuste simplificado{[}cite: 399{]}:

\begin{itemize}
\tightlist
\item
  Convoluciones discretas 2-D (\(N = 2\)),
\item
  Entradas cuadradas (\(i_1 = i_2 = i\)),
\item
  Tamaño de núcleo cuadrado (\(k_1 = k_2 = k\)),
\item
  Mismos pasos a lo largo de ambos ejes (\(s_1 = s_2 = s\)),
\item
  Mismo relleno de ceros a lo largo de ambos ejes (\(p_1 = p_2 = p\)).
\end{itemize}

{[}cite\_start{]}Esto facilita el análisis y la visualización, pero
tenga en cuenta que los resultados aquí descritos también se generalizan
a los casos N-D y no cuadrados{[}cite: 400{]}.

\subsubsection{2.1 Sin relleno de ceros, pasos
unitarios}\label{sin-relleno-de-ceros-pasos-unitarios}

{[}cite\_start{]}El caso más simple de analizar es cuando el núcleo
simplemente se desliza por cada posición de la entrada (es decir,
\(s = 1\) y \(p = 0\)){[}cite: 401{]}. {[}cite\_start{]}La Figura 2.1
proporciona un ejemplo para \(i = 4\) y \(k = 3\){[}cite: 402{]}.

{[}cite\_start{]}Una forma de definir el tamaño de salida en este caso
es por el número de posibles ubicaciones del núcleo en la
entrada{[}cite: 403{]}. {[}cite\_start{]}Consideremos el eje de ancho:
el núcleo comienza en la parte más a la izquierda del mapa de
características de entrada y se desliza en pasos de uno hasta que toca
el lado derecho de la entrada{[}cite: 404{]}. {[}cite\_start{]}El tamaño
de la salida será igual al número de pasos realizados, más uno, lo que
representa la posición inicial del núcleo (Figura 2.8a){[}cite: 404{]}.
{[}cite\_start{]}La misma lógica se aplica al eje de altura{[}cite:
405{]}.

Más formalmente, se puede inferir la siguiente relación:

\textbf{Relación 1.} Para cualquier \(i\) y \(k\), y para \(s = 1\) y
\(p = 0\),

\[o = (i - k) + 1.\]

\begin{center}\rule{0.5\linewidth}{0.5pt}\end{center}

\subsubsection{2.2 Relleno de ceros, pasos
unitarios}\label{relleno-de-ceros-pasos-unitarios}

{[}cite\_start{]}Para tener en cuenta el relleno de ceros (es decir,
solo restringiendo a \(s = 1\)), consideremos su efecto en el tamaño de
entrada efectivo: rellenar con \(p\) ceros cambia el tamaño de entrada
efectivo de \(i\) a \(i + 2p\){[}cite: 407{]}. En el caso general, la
Relación 1 se puede usar para inferir la siguiente relación:

\textbf{Relación 2.} Para cualquier \(i\), \(k\) y \(p\), y para
\(s = 1\),

\[o = (i - k) + 2p + 1.\]

{[}cite\_start{]}La Figura 2.2 proporciona un ejemplo para \(i = 5\),
\(k = 4\) y \(p = 2\){[}cite: 407{]}.

{[}cite\_start{]}En la práctica, se utilizan dos instancias específicas
de relleno de ceros de forma bastante extensa debido a sus respectivas
propiedades{[}cite: 408{]}. Discutámoslas con más detalle.

\paragraph{\texorpdfstring{2.2.1 Relleno a la mitad (\emph{Half (same)
padding})}{2.2.1 Relleno a la mitad (Half (same) padding)}}\label{relleno-a-la-mitad-half-same-padding}

Tener el tamaño de salida igual al tamaño de entrada (es decir,
\(o = i\)) puede ser una propiedad deseable:

\textbf{Relación 3.} Para cualquier \(i\) y para \(k\) impar
(\(k = 2n + 1, n \in \mathbb{N}\)), \(s = 1\) y
\(p = \lfloor k/2 \rfloor = n\),

\[
\begin{aligned}
o &= i + 2\lfloor k/2 \rfloor - (k - 1) \\
&= i + 2n - 2n \\
&= i.
\end{aligned}
\]

{[}cite\_start{]}Esto a veces se denomina \textbf{relleno a la mitad}
(\emph{half}) o \textbf{relleno igual} (\emph{same padding}){[}cite:
409{]}. {[}cite\_start{]}La Figura 2.3 proporciona un ejemplo para
\(i = 5\), \(k = 3\) y (por lo tanto) \(p = 1\){[}cite: 410{]}.

\paragraph{\texorpdfstring{2.2.2 Relleno completo (\emph{Full
padding})}{2.2.2 Relleno completo (Full padding)}}\label{relleno-completo-full-padding}

{[}cite\_start{]}Si bien convolucionar un núcleo generalmente disminuye
el tamaño de salida con respecto al tamaño de entrada, a veces se
requiere lo contrario{[}cite: 411{]}. Esto se puede lograr con un
relleno de ceros adecuado:

\textbf{Relación 4.} Para cualquier \(i\) y \(k\), y para \(p = k - 1\)
y \(s = 1\),

\[
\begin{aligned}
o &= i + 2(k - 1) - (k - 1) \\
&= i + (k - 1).
\end{aligned}
\]

{[}cite\_start{]}Esto a veces se denomina \textbf{relleno completo}
(\emph{full padding}) {[}cite: 418{]}{[}cite\_start{]}, porque en este
ajuste se tiene en cuenta toda posible superposición parcial o completa
del núcleo en el mapa de características de entrada{[}cite: 418{]}.
{[}cite\_start{]}La Figura 2.4 proporciona un ejemplo para \(i = 5\),
\(k = 3\) y (por lo tanto) \(p = 2\){[}cite: 419{]}.

\begin{center}\rule{0.5\linewidth}{0.5pt}\end{center}

{[}cite\_start{]}\textbf{Figura 2.1:} (Sin relleno, pasos unitarios)
Convolución de un núcleo de \(3 \times 3\) sobre una entrada de
\(4 \times 4\) utilizando pasos unitarios (es decir,
\(i = 4, k = 3, s = 1\) y \(p = 0\)){[}cite: 413{]}.

{[}cite\_start{]}\textbf{Figura 2.2:} (Relleno arbitrario, pasos
unitarios) Convolución de un núcleo de \(4 \times 4\) sobre una entrada
de \(5 \times 5\) rellenada con un borde de ceros de \(2 \times 2\)
utilizando pasos unitarios (es decir, \(i = 5, k = 4, s = 1\) y
\(p = 2\)){[}cite: 414{]}.

{[}cite\_start{]}\textbf{Figura 2.3:} (Relleno a la mitad, pasos
unitarios) Convolución de un núcleo de \(3 \times 3\) sobre una entrada
de \(5 \times 5\) utilizando relleno a la mitad y pasos unitarios (es
decir, \(i = 5, k = 3, s = 1\) y \(p = 1\)){[}cite: 415{]}.

{[}cite\_start{]}\textbf{Figura 2.4:} (Relleno completo, pasos
unitarios) Convolución de un núcleo de \(3 \times 3\) sobre una entrada
de \(5 \times 5\) utilizando relleno completo y pasos unitarios (es
decir, \(i = 5, k = 3, s = 1\) y \(p = 2\)){[}cite: 416{]}.

\begin{center}\rule{0.5\linewidth}{0.5pt}\end{center}

\subsubsection{2.3 Sin relleno de ceros, pasos no
unitarios}\label{sin-relleno-de-ceros-pasos-no-unitarios}

{[}cite\_start{]}Todas las relaciones derivadas hasta ahora solo se
aplican a convoluciones con pasos unitarios{[}cite: 420{]}.
{[}cite\_start{]}Incorporar \textbf{pasos no unitarios} requiere otro
salto de inferencia{[}cite: 421{]}. {[}cite\_start{]}Para facilitar el
análisis, ignoremos momentáneamente el relleno de ceros (es decir,
\(s > 1\) y \(p = 0\)){[}cite: 422{]}. {[}cite\_start{]}La Figura 2.5
proporciona un ejemplo para \(i = 5\), \(k = 3\) y \(s = 2\){[}cite:
422{]}.

{[}cite\_start{]}Una vez más, el tamaño de la salida se puede definir en
términos del número de posibles ubicaciones del núcleo en la
entrada{[}cite: 423{]}. {[}cite\_start{]}Consideremos el eje de ancho:
el núcleo comienza como de costumbre en la parte más a la izquierda de
la entrada, pero esta vez se desliza en pasos de tamaño \(s\) hasta que
toca el lado derecho de la entrada{[}cite: 424{]}. {[}cite\_start{]}El
tamaño de la salida es de nuevo igual al número de pasos realizados, más
uno, lo que representa la posición inicial del núcleo (Figura
2.8b){[}cite: 425{]}. {[}cite\_start{]}La misma lógica se aplica al eje
de altura{[}cite: 425{]}.

A partir de esto, se puede inferir la siguiente relación:

\textbf{Relación 5.} Para cualquier \(i\), \(k\) y \(s\), y para
\(p = 0\),

\[o = \left\lfloor \frac{i - k}{s} \right\rfloor + 1.\]

{[}cite\_start{]}La función piso (\(\lfloor \cdot \rfloor\)) da cuenta
del hecho de que a veces el último paso posible no coincide con que el
núcleo llegue al final de la entrada, es decir, algunas unidades de
entrada quedan fuera (consulte la Figura 2.7 para un ejemplo de tal
caso){[}cite: 426{]}.

\subsubsection{2.4 Relleno de ceros, pasos no
unitarios}\label{relleno-de-ceros-pasos-no-unitarios}

El caso más general (convolucionar sobre una entrada rellenada con ceros
utilizando pasos no unitarios) se puede derivar aplicando la Relación 5
a una entrada efectiva de tamaño \(i + 2p\), en analogía con lo que se
hizo para la Relación 2:

\textbf{Relación 6.} Para cualquier \(i\), \(k\), \(p\) y \(s\),

\[o = \left\lfloor \frac{i + 2p - k}{s} \right\rfloor + 1.\]

{[}cite\_start{]}Como antes, la función piso significa que en algunos
casos una convolución producirá el mismo tamaño de salida para múltiples
tamaños de entrada{[}cite: 427{]}. {[}cite\_start{]}Más específicamente,
si \(i + 2p - k\) es un múltiplo de \(s\), entonces cualquier tamaño de
entrada \(j = i + a\), \(a \in \{0, \dots, s-1\}\) producirá el mismo
tamaño de salida{[}cite: 428{]}. {[}cite\_start{]}Tenga en cuenta que
esta ambigüedad se aplica solo para \(s > 1\){[}cite: 429{]}.

{[}cite\_start{]}La Figura 2.6 muestra un ejemplo con
\(i = 5, k = 3, s = 2\) y \(p = 1\), mientras que la Figura 2.7
proporciona un ejemplo para \(i = 6, k = 3, s = 2\) y \(p = 1\){[}cite:
429{]}. {[}cite\_start{]}Curiosamente, a pesar de tener diferentes
tamaños de entrada, estas convoluciones comparten el mismo tamaño de
salida{[}cite: 430{]}. {[}cite\_start{]}Si bien esto no afecta el
análisis de las convoluciones, complicará el análisis en el caso de las
convoluciones transpuestas{[}cite: 431{]}.

\begin{center}\rule{0.5\linewidth}{0.5pt}\end{center}

{[}cite\_start{]}\textbf{Figura 2.5:} (Sin relleno de ceros, pasos
arbitrarios) Convolución de un núcleo de \(3 \times 3\) sobre una
entrada de \(5 \times 5\) utilizando pasos de \(2 \times 2\) (es decir,
\(i = 5, k = 3, s = 2\) y \(p = 0\)){[}cite: 432{]}.

{[}cite\_start{]}\textbf{Figura 2.6:} (Relleno y pasos arbitrarios)
Convolución de un núcleo de \(3 \times 3\) sobre una entrada de
\(5 \times 5\) rellenada con un borde de ceros de \(1 \times 1\)
utilizando pasos de \(2 \times 2\) (es decir, \(i = 5, k = 3, s = 2\) y
\(p = 1\)){[}cite: 433{]}.

{[}cite\_start{]}\textbf{Figura 2.7:} (Relleno y pasos arbitrarios)
Convolución de un núcleo de \(3 \times 3\) sobre una entrada de
\(6 \times 6\) rellenada con un borde de ceros de \(1 \times 1\)
utilizando pasos de \(2 \times 2\) (es decir, \(i = 6, k = 3, s = 2\) y
\(p = 1\)){[}cite: 434{]}. {[}cite\_start{]}En este caso, la fila
inferior y la columna derecha de la entrada rellenada con ceros no están
cubiertas por el núcleo{[}cite: 435{]}.

{[}cite\_start{]}\textbf{Figura 2.8:} Conteo de posiciones del
núcleo{[}cite: 438{]}. (a) {[}cite\_start{]}El núcleo tiene que
deslizarse dos pasos hacia la derecha para tocar el lado derecho de la
entrada (y de manera equivalente hacia abajo){[}cite: 436{]}.
{[}cite\_start{]}Añadiendo uno para tener en cuenta la posición inicial
del núcleo, el tamaño de salida es \(3 \times 3\){[}cite: 437{]}. (b)
{[}cite\_start{]}El núcleo tiene que deslizarse un paso de tamaño dos
hacia la derecha para tocar el lado derecho de la entrada (y de manera
equivalente hacia abajo){[}cite: 437{]}. {[}cite\_start{]}Añadiendo uno
para tener en cuenta la posición inicial del núcleo, el tamaño de salida
es \(2 \times 2\){[}cite: 438{]}.

\begin{center}\rule{0.5\linewidth}{0.5pt}\end{center}

\subsection{Capítulo 3 Aritmética del
Pooling}\label{capuxedtulo-3-aritmuxe9tica-del-pooling}

{[}cite\_start{]}En una red neuronal, las capas de \emph{pooling}
proporcionan \textbf{invariancia a pequeñas traslaciones de la
entrada}{[}cite: 440{]}. {[}cite\_start{]}El tipo más común de
\emph{pooling} es el \textbf{max pooling} (agrupamiento máximo), que
consiste en dividir la entrada en parches (generalmente no superpuestos)
y generar el valor máximo de cada parche{[}cite: 441{]}.
{[}cite\_start{]}Existen otros tipos de \emph{pooling}, por ejemplo, el
\emph{mean} o \emph{average pooling} (agrupamiento medio o promedio),
que comparten la misma idea de agregar la entrada localmente aplicando
una no linealidad al contenido de algunos parches (Boureau et al.,
2010a,b, 2011; Saxe et al., 2011){[}cite: 442{]}.

{[}cite\_start{]}Algunos lectores pueden haber notado que el tratamiento
de la aritmética de la convolución solo se basa en la suposición de que
alguna función se aplica repetidamente a subconjuntos de la
entrada{[}cite: 444{]}. {[}cite\_start{]}Esto significa que las
relaciones derivadas en el capítulo anterior se pueden reutilizar en el
caso de la aritmética de \emph{pooling}{[}cite: 444{]}. Dado que el
\emph{pooling} no implica relleno de ceros, la relación que describe el
caso general es la siguiente:

\textbf{Relación 7.} Para cualquier \(i\), \(k\) y \(s\),

\[o = \left\lfloor \frac{i - k}{s} \right\rfloor + 1.\]

{[}cite\_start{]}Esta relación es válida para cualquier tipo de
\emph{pooling}{[}cite: 446{]}.

\begin{center}\rule{0.5\linewidth}{0.5pt}\end{center}

\subsection{Capítulo 4 Aritmética de la convolución
transpuesta}\label{capuxedtulo-4-aritmuxe9tica-de-la-convoluciuxf3n-transpuesta}

{[}cite\_start{]}La necesidad de \textbf{convoluciones transpuestas}
generalmente surge del deseo de utilizar una transformación que vaya en
la dirección opuesta a una convolución normal {[}cite:
447{]}{[}cite\_start{]}, es decir, de algo que tenga la forma de la
salida de alguna convolución a algo que tenga la forma de su entrada,
manteniendo un patrón de conectividad que sea compatible con dicha
convolución{[}cite: 447{]}. {[}cite\_start{]}Por ejemplo, se podría usar
tal transformación como la capa de decodificación de un autoencoder
convolucional o para proyectar mapas de características a un espacio de
mayor dimensión{[}cite: 448{]}.

{[}cite\_start{]}Una vez más, el caso convolucional es considerablemente
más complejo que el caso totalmente conectado, que solo requiere usar
una matriz de pesos cuya forma se ha traspuesto{[}cite: 449{]}.
{[}cite\_start{]}Sin embargo, dado que toda convolución se reduce a una
implementación eficiente de una operación matricial, los conocimientos
obtenidos del caso totalmente conectado son útiles para resolver el caso
convolucional{[}cite: 450{]}.

{[}cite\_start{]}Al igual que con la aritmética de la convolución, la
disertación sobre la aritmética de la convolución transpuesta se
simplifica por el hecho de que las propiedades de la convolución
transpuesta \textbf{no interactúan a través de los ejes}{[}cite: 452{]}.
{[}cite\_start{]}El capítulo se centrará en el siguiente ajuste{[}cite:
452{]}:

\begin{itemize}
\tightlist
\item
  Convoluciones transpuestas 2-D (\(N = 2\)),
\item
  Entradas cuadradas (\(i_1 = i_2 = i\)),
\item
  Tamaño de núcleo cuadrado (\(k_1 = k_2 = k\)),
\item
  Mismos pasos a lo largo de ambos ejes (\(s_1 = s_2 = s\)),
\item
  Mismo relleno de ceros a lo largo de ambos ejes (\(p_1 = p_2 = p\)).
\end{itemize}

{[}cite\_start{]}Una vez más, los resultados descritos se generalizan a
los casos N-D y no cuadrados{[}cite: 453{]}.

\subsubsection{4.1 La convolución como una operación
matricial}\label{la-convoluciuxf3n-como-una-operaciuxf3n-matricial}

{[}cite\_start{]}Tomemos como ejemplo la convolución representada en la
Figura 2.1{[}cite: 454{]}. {[}cite\_start{]}Si la entrada y la salida se
desenrollaran en vectores de izquierda a derecha, de arriba a abajo, la
convolución podría representarse como una \textbf{matriz dispersa \(C\)}
donde los elementos no nulos son los elementos \(w_{i,j}\) del núcleo
(siendo \(i\) y \(j\) la fila y la columna del núcleo,
respectivamente){[}cite: 455{]}:

\[
\begin{pmatrix}
w_{0,0} & w_{0,1} & w_{0,2} & 0 & w_{1,0} & w_{1,1} & w_{1,2} & 0 & w_{2,0} & w_{2,1} & w_{2,2} & 0 & 0 & 0 & 0 & 0 \\
0 & w_{0,0} & w_{0,1} & w_{0,2} & 0 & w_{1,0} & w_{1,1} & w_{1,2} & 0 & w_{2,0} & w_{2,1} & w_{2,2} & 0 & 0 & 0 & 0 \\
0 & 0 & 0 & 0 & w_{0,0} & w_{0,1} & w_{0,2} & 0 & w_{1,0} & w_{1,1} & w_{1,2} & 0 & w_{2,0} & w_{2,1} & w_{2,2} & 0 \\
0 & 0 & 0 & 0 & 0 & w_{0,0} & w_{0,1} & w_{0,2} & 0 & w_{1,0} & w_{1,1} & w_{1,2} & 0 & w_{2,0} & w_{2,1} & w_{2,2}
\end{pmatrix}
\]

{[}cite\_start{]}Esta operación lineal toma la matriz de entrada
aplanada como un vector de 16 dimensiones y produce un vector de 4
dimensiones que luego se remodela como la matriz de salida de
\(2 \times 2\){[}cite: 457{]}.

{[}cite\_start{]}Usando esta representación, el \textbf{pase hacia
atrás} (\emph{backward pass}) se obtiene fácilmente transponiendo \(C\)
{[}cite: 458{]}{[}cite\_start{]}; en otras palabras, el error se
retropropaga multiplicando la pérdida por \(C^T\){[}cite: 459{]}.
{[}cite\_start{]}Esta operación toma un vector de 4 dimensiones como
entrada y produce un vector de 16 dimensiones como salida, y su patrón
de conectividad es compatible con \(C\) por construcción{[}cite: 460{]}.

{[}cite\_start{]}Notablemente, el núcleo \(w\) define tanto las matrices
\(C\) como \(C^T\) utilizadas para los pases hacia adelante y hacia
atrás{[}cite: 461{]}.

\subsubsection{4.2 Convolución
transpuesta}\label{convoluciuxf3n-transpuesta}

{[}cite\_start{]}Consideremos ahora lo que se requeriría para ir en la
otra dirección, es decir, mapear de un espacio de 4 dimensiones a un
espacio de 16 dimensiones, manteniendo el patrón de conectividad de la
convolución representada en la Figura 2.1{[}cite: 462{]}.
{[}cite\_start{]}Esta operación se conoce como \textbf{convolución
transpuesta}{[}cite: 462{]}.

{[}cite\_start{]}Las convoluciones transpuestas, también llamadas
\textbf{convoluciones con pasos fraccionados} o
\textbf{deconvoluciones}\textsuperscript{1}, funcionan intercambiando
los pases hacia adelante y hacia atrás de una convolución{[}cite:
463{]}. {[}cite\_start{]}Una forma de expresarlo es notar que el núcleo
define una convolución, pero si es una convolución directa o una
convolución transpuesta está determinado por cómo se calculan los pases
hacia adelante y hacia atrás{[}cite: 464{]}. {[}cite\_start{]}Por
ejemplo, aunque el núcleo \(w\) define una convolución cuyos pases hacia
adelante y hacia atrás se calculan multiplicando por \(C\) y \(C^T\)
respectivamente, también define una convolución transpuesta cuyos pases
hacia adelante y hacia atrás se calculan multiplicando por \(C^T\) y
\((C^T)^T = C\) respectivamente\textsuperscript{2}{[}cite: 464{]}.

{[}cite\_start{]}Finalmente, tenga en cuenta que \textbf{siempre es
posible emular una convolución transpuesta con una convolución
directa}{[}cite: 465{]}. {[}cite\_start{]}La desventaja es que
generalmente implica agregar muchas columnas y filas de ceros a la
entrada, lo que resulta en una implementación mucho menos
eficiente{[}cite: 468{]}.

{[}cite\_start{]}Basándose en lo que se ha introducido hasta ahora, este
capítulo procederá de forma algo inversa con respecto al capítulo de
aritmética de la convolución, derivando las propiedades de cada
convolución transpuesta haciendo referencia a la convolución directa con
la que comparte el núcleo, y definiendo la convolución directa
equivalente{[}cite: 469{]}.

{[}cite\_start{]}\textsuperscript{1}El término
``\textbf{deconvolución}'' se utiliza a veces en la literatura, pero
abogamos en contra de él basándonos en que una deconvolución se define
matemáticamente como la inversa de una convolución, que es diferente de
una convolución transpuesta{[}cite: 466{]}.

{[}cite\_start{]}\textsuperscript{2}La operación de convolución
transpuesta puede considerarse como el gradiente de alguna convolución
con respecto a su entrada, que es generalmente cómo se implementan las
convoluciones transpuestas en la práctica{[}cite: 467{]}.

\begin{center}\rule{0.5\linewidth}{0.5pt}\end{center}

\subsubsection{4.3 Sin relleno de ceros, pasos unitarios,
traspuesta}\label{sin-relleno-de-ceros-pasos-unitarios-traspuesta}

{[}cite\_start{]}La forma más sencilla de pensar en una convolución
transpuesta en una entrada dada es imaginar dicha entrada como el
resultado de una convolución directa aplicada a algún mapa de
características inicial{[}cite: 471{]}. {[}cite\_start{]}La convolución
transpuesta puede considerarse entonces como la operación que permite
recuperar la forma\textsuperscript{3} de este mapa de características
inicial{[}cite: 472{]}.

{[}cite\_start{]}Consideremos la convolución de un núcleo de
\(3 \times 3\) en una entrada de \(4 \times 4\) con paso unitario y sin
relleno (es decir, \(i = 4, k = 3, s = 1\) y \(p = 0\)){[}cite: 473{]}.
{[}cite\_start{]}Como se representa en la Figura 2.1, esto produce una
salida de \(2 \times 2\){[}cite: 474{]}. {[}cite\_start{]}La transpuesta
de esta convolución tendrá entonces una salida con forma de
\(4 \times 4\) cuando se aplique a una entrada de \(2 \times 2\){[}cite:
475{]}.

{[}cite\_start{]}Otra forma de obtener el resultado de una convolución
transpuesta es aplicar una convolución directa equivalente, pero mucho
menos eficiente{[}cite: 476{]}. {[}cite\_start{]}El ejemplo descrito
hasta ahora podría abordarse convolucionando un núcleo de \(3 \times 3\)
sobre una entrada de \(2 \times 2\) rellenada con un borde de ceros de
\(2 \times 2\) utilizando pasos unitarios (es decir,
\(i' = 2, k' = k, s' = 1\) y \(p' = 2\)), como se muestra en la Figura
4.1{[}cite: 477{]}. {[}cite\_start{]}Notablemente, los tamaños del
núcleo y del paso siguen siendo los mismos, pero la entrada de la
convolución transpuesta ahora está rellenada con
ceros\textsuperscript{4}{[}cite: 477{]}.

{[}cite\_start{]}Una forma de comprender la lógica detrás del relleno de
ceros es considerar el patrón de conectividad de la convolución
transpuesta y usarlo para guiar el diseño de la convolución
equivalente{[}cite: 478{]}. {[}cite\_start{]}Por ejemplo, el píxel
superior izquierdo de la entrada de la convolución directa solo
contribuye al píxel superior izquierdo de la salida, el píxel superior
derecho solo está conectado al píxel de salida superior derecho, y así
sucesivamente{[}cite: 479{]}. {[}cite\_start{]}Para mantener el mismo
patrón de conectividad en la convolución equivalente, es necesario
rellenar la entrada con ceros de tal manera que la primera aplicación
(superior izquierda) del núcleo solo toque el píxel superior izquierdo,
es decir, el relleno tiene que ser igual al tamaño del núcleo menos
uno{[}cite: 480{]}.

Procediendo de la misma manera, es posible determinar observaciones
similares para los otros elementos de la imagen, dando lugar a la
siguiente relación:

\textbf{Relación 8.} Una convolución descrita por \(s = 1, p = 0\) y
\(k\) tiene una convolución transpuesta asociada descrita por
\(k' = k, s' = s\) y \(p' = k - 1\) y su tamaño de salida es

\[o' = i' + (k - 1).\]

{[}cite\_start{]}Curiosamente, esto corresponde a una convolución
completamente rellenada con pasos unitarios{[}cite: 485{]}.

{[}cite\_start{]}\textsuperscript{3}Tenga en cuenta que la convolución
transpuesta no garantiza que recupere la entrada en sí, ya que no se
define como la inversa de la convolución, sino que simplemente devuelve
un mapa de características que tiene el mismo ancho y alto{[}cite:
480{]}.

{[}cite\_start{]}\textsuperscript{4}Tenga en cuenta que, aunque es
equivalente a aplicar la matriz transpuesta, esta visualización agrega
muchas multiplicaciones por cero en forma de relleno de ceros{[}cite:
481{]}. {[}cite\_start{]}Esto se hace aquí con fines ilustrativos, pero
es ineficiente, y las implementaciones de \emph{software} normalmente no
realizarán las multiplicaciones inútiles por cero{[}cite: 482{]}.

\begin{center}\rule{0.5\linewidth}{0.5pt}\end{center}

\subsubsection{4.4 Relleno de ceros, pasos unitarios,
traspuesta}\label{relleno-de-ceros-pasos-unitarios-traspuesta}

{[}cite\_start{]}Sabiendo que la transpuesta de una convolución sin
relleno es equivalente a convolucionar una entrada rellenada con ceros,
sería razonable suponer que la transpuesta de una convolución rellenada
con ceros es equivalente a convolucionar una entrada rellenada con menos
ceros{[}cite: 487{]}. {[}cite\_start{]}De hecho, es el caso, como se
muestra en la Figura 4.2 para \(i = 5, k = 4\) y \(p = 2\){[}cite:
487{]}.

Formalmente, la siguiente relación se aplica a las convoluciones
rellenadas con ceros:

\textbf{Relación 9.} Una convolución descrita por \(s = 1, k\) y \(p\)
tiene una convolución transpuesta asociada descrita por
\(k' = k, s' = s\) y \(p' = k - p - 1\) y su tamaño de salida es

\[o' = i' + (k - 1) - 2p.\]

\paragraph{\texorpdfstring{4.4.1 Relleno a la mitad (\emph{Half (same)
padding}),
traspuesta}{4.4.1 Relleno a la mitad (Half (same) padding), traspuesta}}\label{relleno-a-la-mitad-half-same-padding-traspuesta}

{[}cite\_start{]}Aplicando el mismo razonamiento inductivo que antes, es
razonable esperar que la convolución equivalente de la transpuesta de
una convolución con relleno a la mitad sea ella misma una convolución
con relleno a la mitad, dado que el tamaño de salida de una convolución
con relleno a la mitad es el mismo que su tamaño de entrada{[}cite:
489{]}. Por lo tanto, se aplica la siguiente relación:

\textbf{Relación 10.} Una convolución descrita por
\(k = 2n + 1, n \in \mathbb{N}, s = 1\) y
\(p = \lfloor k/2 \rfloor = n\) tiene una convolución transpuesta
asociada descrita por \(k' = k, s' = s\) y \(p' = p\) y su tamaño de
salida es

\[
\begin{aligned}
o' &= i' + (k - 1) - 2p \\
&= i' + 2n - 2n \\
&= i'.
\end{aligned}
\]

{[}cite\_start{]}La Figura 4.3 proporciona un ejemplo para
\(i = 5, k = 3\) y (por lo tanto) \(p = 1\){[}cite: 490{]}.

\paragraph{\texorpdfstring{4.4.2 Relleno completo (\emph{Full padding}),
traspuesta}{4.4.2 Relleno completo (Full padding), traspuesta}}\label{relleno-completo-full-padding-traspuesta}

Sabiendo que la convolución equivalente de la transpuesta de una
convolución sin relleno implica relleno completo, no es sorprendente que
la equivalente de la transpuesta de una convolución completamente
rellenada sea una convolución sin relleno:

\textbf{Relación 11.} Una convolución descrita por \(s = 1, k\) y
\(p = k - 1\) tiene una convolución transpuesta asociada descrita por
\(k' = k, s' = s\) y \(p' = 0\) y su tamaño de salida es

\[
\begin{aligned}
o' &= i' + (k - 1) - 2p \\
&= i' - (k - 1)
\end{aligned}
\]

{[}cite\_start{]}La Figura 4.4 proporciona un ejemplo para
\(i = 5, k = 3\) y (por lo tanto) \(p = 2\){[}cite: 498{]}.

\begin{center}\rule{0.5\linewidth}{0.5pt}\end{center}

{[}cite\_start{]}\textbf{Figura 4.1:} La transpuesta de convolucionar un
núcleo de \(3 \times 3\) sobre una entrada de \(4 \times 4\) utilizando
pasos unitarios (es decir, \(i = 4, k = 3, s = 1\) y \(p = 0\)){[}cite:
491{]}. {[}cite\_start{]}Es equivalente a convolucionar un núcleo de
\(3 \times 3\) sobre una entrada de \(2 \times 2\) rellenada con un
borde de ceros de \(2 \times 2\) utilizando pasos unitarios (es decir,
\(i' = 2, k' = k, s' = 1\) y \(p' = 2\)){[}cite: 492{]}.

{[}cite\_start{]}\textbf{Figura 4.2:} La transpuesta de convolucionar un
núcleo de \(4 \times 4\) sobre una entrada de \(5 \times 5\) rellenada
con un borde de ceros de \(2 \times 2\) utilizando pasos unitarios (es
decir, \(i = 5, k = 4, s = 1\) y \(p = 2\)){[}cite: 493{]}.
{[}cite\_start{]}Es equivalente a convolucionar un núcleo de
\(4 \times 4\) sobre una entrada de \(6 \times 6\) rellenada con un
borde de ceros de \(1 \times 1\) utilizando pasos unitarios (es decir,
\(i' = 6, k' = k, s' = 1\) y \(p' = 1\)){[}cite: 494{]}.

{[}cite\_start{]}\textbf{Figura 4.3:} La transpuesta de convolucionar un
núcleo de \(3 \times 3\) sobre una entrada de \(5 \times 5\) utilizando
relleno a la mitad y pasos unitarios (es decir, \(i = 5, k = 3, s = 1\)
y \(p = 1\)){[}cite: 495{]}. {[}cite\_start{]}Es equivalente a
convolucionar un núcleo de \(3 \times 3\) sobre una entrada de
\(5 \times 5\) utilizando relleno a la mitad y pasos unitarios (es
decir, \(i' = 5, k' = k, s' = 1\) y \(p' = 1\)){[}cite: 496{]}.

\begin{center}\rule{0.5\linewidth}{0.5pt}\end{center}

\subsubsection{4.5 Sin relleno de ceros, pasos no unitarios,
traspuesta}\label{sin-relleno-de-ceros-pasos-no-unitarios-traspuesta}

{[}cite\_start{]}Utilizando el mismo tipo de lógica inductiva que para
las convoluciones rellenadas con ceros, se podría esperar que la
transpuesta de una convolución con \(s > 1\) involucre una convolución
equivalente con \(s < 1\){[}cite: 499{]}. {[}cite\_start{]}Como se
explicará, esta es una intuición válida, por lo que las convoluciones
transpuestas a veces se denominan \textbf{convoluciones con pasos
fraccionados}{[}cite: 499{]}.

{[}cite\_start{]}La Figura 4.5 proporciona un ejemplo para
\(i = 5, k = 3\) y \(s = 2\) que ayuda a comprender lo que implican los
pasos fraccionados: \textbf{se insertan ceros entre las unidades de
entrada}, lo que hace que el núcleo se mueva a un ritmo más lento que
con los pasos unitarios\textsuperscript{5}{[}cite: 500{]}.
{[}cite\_start{]}Por el momento, se asumirá que la convolución no está
rellenada (\(p = 0\)) y que su tamaño de entrada \(i\) es tal que
\(i - k\) es un múltiplo de \(s\){[}cite: 501{]}. En ese caso, se cumple
la siguiente relación:

\textbf{Relación 12.} Una convolución descrita por \(p = 0, k\) y \(s\)
y cuyo tamaño de entrada es tal que \(i - k\) es un múltiplo de \(s\),
tiene una convolución transpuesta asociada descrita por
\(\tilde{i}', k' = k, s' = 1\) y \(p' = k - 1\), donde \(\tilde{i}'\) es
el tamaño de la entrada estirada obtenida al agregar \(s - 1\) ceros
entre cada unidad de entrada, y su tamaño de salida es

\[o' = s(i' - 1) + k.\]

{[}cite\_start{]}\textsuperscript{5}Hacer esto es ineficiente y las
implementaciones en el mundo real evitan multiplicaciones inútiles por
cero, pero conceptualmente es cómo se puede pensar en la transpuesta de
una convolución con pasos{[}cite: 503{]}.

\begin{center}\rule{0.5\linewidth}{0.5pt}\end{center}

\subsubsection{4.6 Relleno de ceros, pasos no unitarios,
traspuesta}\label{relleno-de-ceros-pasos-no-unitarios-traspuesta}

Cuando el tamaño de entrada de la convolución \(i\) es tal que
\(i + 2p - k\) es un múltiplo de \(s\), el análisis se puede extender al
caso rellenado con ceros combinando la Relación 9 y la Relación 12:

\textbf{Relación 13.} Una convolución descrita por \(k, s\) y \(p\) y
cuyo tamaño de entrada \(i\) es tal que \(i + 2p - k\) es un múltiplo de
\(s\) tiene una convolución transpuesta asociada descrita por
\(\tilde{i}', k' = k, s' = 1\) y \(p' = k - p - 1\), donde
\(\tilde{i}'\) es el tamaño de la entrada estirada obtenida al agregar
\(s - 1\) ceros entre cada unidad de entrada, y su tamaño de salida es

\[o' = s(i' - 1) + k - 2p.\]

{[}cite\_start{]}La Figura 4.6 proporciona un ejemplo para
\(i = 5, k = 3, s = 2\) y \(p = 1\){[}cite: 512{]}.

{[}cite\_start{]}La restricción en el tamaño de la entrada \(i\) se
puede relajar introduciendo otro parámetro \(a \in \{0, \dots, s - 1\}\)
que permite distinguir entre los \(s\) casos diferentes que conducen al
mismo \(i'\){[}cite: 513{]}:

\textbf{Relación 14.} Una convolución descrita por \(k, s\) y \(p\)
tiene una convolución transpuesta asociada descrita por
\(a, \tilde{i}', k' = k, s' = 1\) y \(p' = k - p - 1\), donde
\(\tilde{i}'\) es el tamaño de la entrada estirada obtenida al agregar
\(s - 1\) ceros entre cada unidad de entrada, y
\(a = (i + 2p - k) \mod s\) representa el número de ceros agregados a
los bordes inferior y derecho de la entrada, y su tamaño de salida es

\[o' = s(i' - 1) + a + k - 2p.\]

{[}cite\_start{]}La Figura 4.7 proporciona un ejemplo para
\(i = 6, k = 3, s = 2\) y \(p = 1\){[}cite: 515{]}.

\begin{center}\rule{0.5\linewidth}{0.5pt}\end{center}

{[}cite\_start{]}\textbf{Figura 4.4:} La transpuesta de convolucionar un
núcleo de \(3 \times 3\) sobre una entrada de \(5 \times 5\) utilizando
relleno completo y pasos unitarios (es decir, \(i = 5, k = 3, s = 1\) y
\(p = 2\)){[}cite: 504{]}. {[}cite\_start{]}Es equivalente a
convolucionar un núcleo de \(3 \times 3\) sobre una entrada de
\(7 \times 7\) utilizando pasos unitarios (es decir,
\(i' = 7, k' = k, s' = 1\) y \(p' = 0\)){[}cite: 505{]}.

{[}cite\_start{]}\textbf{Figura 4.5:} La transpuesta de convolucionar un
núcleo de \(3 \times 3\) sobre una entrada de \(5 \times 5\) utilizando
pasos de \(2 \times 2\) (es decir, \(i = 5, k = 3, s = 2\) y
\(p = 0\)){[}cite: 506{]}. {[}cite\_start{]}Es equivalente a
convolucionar un núcleo de \(3 \times 3\) sobre una entrada de
\(2 \times 2\) (con 1 cero insertado entre entradas) rellenada con un
borde de ceros de \(2 \times 2\) utilizando pasos unitarios (es decir,
\(i' = 2, \tilde{i}' = 3, k' = k, s' = 1\) y \(p' = 2\)){[}cite: 507{]}.

{[}cite\_start{]}\textbf{Figura 4.6:} La transpuesta de convolucionar un
núcleo de \(3 \times 3\) sobre una entrada de \(5 \times 5\) rellenada
con un borde de ceros de \(1 \times 1\) utilizando pasos de
\(2 \times 2\) (es decir, \(i = 5, k = 3, s = 2\) y \(p = 1\)){[}cite:
508{]}. {[}cite\_start{]}Es equivalente a convolucionar un núcleo de
\(3 \times 3\) sobre una entrada de \(3 \times 3\) (con 1 cero insertado
entre entradas) rellenada con un borde de ceros de \(1 \times 1\)
utilizando pasos unitarios (es decir,
\(i' = 3, \tilde{i}' = 5, k' = k, s' = 1\) y \(p' = 1\)){[}cite: 509{]}.

{[}cite\_start{]}\textbf{Figura 4.7:} La transpuesta de convolucionar un
núcleo de \(3 \times 3\) sobre una entrada de \(6 \times 6\) rellenada
con un borde de ceros de \(1 \times 1\) utilizando pasos de
\(2 \times 2\) (es decir, \(i = 6, k = 3, s = 2\) y \(p = 1\)){[}cite:
516{]}. {[}cite\_start{]}Es equivalente a convolucionar un núcleo de
\(3 \times 3\) sobre una entrada de \(2 \times 2\) (con 1 cero insertado
entre entradas) rellenada con un borde de ceros de \(1 \times 1\) (con
un borde adicional de tamaño 1 agregado a los bordes inferior y derecho)
utilizando pasos unitarios (es decir,
\(i' = 3, \tilde{i}' = 5, a = 1, k' = k, s' = 1\) y \(p' = 1\)){[}cite:
517{]}.

\begin{center}\rule{0.5\linewidth}{0.5pt}\end{center}

\subsection{Capítulo 5 Convoluciones
misceláneas}\label{capuxedtulo-5-convoluciones-misceluxe1neas}

\subsubsection{5.1 Convoluciones
dilatadas}\label{convoluciones-dilatadas}

{[}cite\_start{]}Los lectores familiarizados con la literatura de
aprendizaje profundo pueden haber notado el término
``\textbf{convoluciones dilatadas}'' (o ``\textbf{convoluciones \emph{à
trous}}'', de la expresión francesa \emph{convolutions à trous})
aparecer en artículos recientes{[}cite: 519{]}. {[}cite\_start{]}Aquí
intentamos proporcionar una comprensión intuitiva de las convoluciones
dilatadas{[}cite: 520{]}. Para una descripción más profunda y para
comprender en qué contextos se aplican, consulte Chen et al.~(2014);
{[}cite\_start{]}Yu y Koltun (2015){[}cite: 522{]}.

{[}cite\_start{]}Las convoluciones dilatadas ``\textbf{inflan}'' el
núcleo insertando espacios entre los elementos del núcleo{[}cite:
523{]}. {[}cite\_start{]}La ``\textbf{tasa}'' de dilatación está
controlada por un hiperparámetro adicional \(d\){[}cite: 524{]}.
{[}cite\_start{]}Las implementaciones pueden variar, pero generalmente
se insertan \(d - 1\) espacios entre los elementos del núcleo, de modo
que \(d = 1\) corresponde a una convolución regular{[}cite: 525{]}.

{[}cite\_start{]}Las convoluciones dilatadas se utilizan para
\textbf{aumentar de forma económica el campo receptivo} de las unidades
de salida sin aumentar el tamaño del núcleo {[}cite:
526{]}{[}cite\_start{]}, lo que es especialmente efectivo cuando se
apilan múltiples convoluciones dilatadas una tras otra{[}cite: 526{]}.
Para un ejemplo concreto, consulte Oord et al.~(2016){[}cite\_start{]},
en el que el modelo WaveNet propuesto implementa un modelo generativo
autorregresivo para audio sin procesar que utiliza convoluciones
dilatadas para \textbf{condicionar nuevos fotogramas de audio en un gran
contexto de fotogramas de audio pasados}{[}cite: 527{]}.

{[}cite\_start{]}Para comprender la relación que une la tasa de
dilatación \(d\) y el tamaño de salida \(o\), es útil pensar en el
impacto de \(d\) en el \textbf{tamaño efectivo del núcleo}{[}cite:
528{]}. Un núcleo de tamaño \(k\) dilatado por un factor \(d\) tiene un
tamaño efectivo

\[\hat{k} = k + (k - 1)(d - 1).\]

Esto se puede combinar con la Relación 6 para formar la siguiente
relación para las convoluciones dilatadas:

\textbf{Relación 15.} Para cualquier \(i, k, p\) y \(s\), y para una
tasa de dilatación \(d\),

\[o = \left\lfloor \frac{i + 2p - k - (k - 1)(d - 1)}{s} \right\rfloor + 1.\]

{[}cite\_start{]}La Figura 5.1 proporciona un ejemplo para
\(i = 7, k = 3\) y \(d = 2\){[}cite: 531{]}.

\begin{center}\rule{0.5\linewidth}{0.5pt}\end{center}

{[}cite\_start{]}\textbf{Figura 5.1:} (Convolución dilatada)
Convolucionar un núcleo de \(3 \times 3\) sobre una entrada de
\(7 \times 7\) con un factor de dilatación de 2 (es decir,
\(i = 7, k = 3, d = 2, s = 1\) y \(p = 0\)){[}cite: 530{]}.

\begin{center}\rule{0.5\linewidth}{0.5pt}\end{center}

\section{Bibliography styles}\label{bibliography-styles}

Here are two sample references: \citet{Feynman1963118}
\citet{Dirac1953888}.

By default, natbib will be used with the \texttt{authoryear} style, set
in \texttt{classoption} variable in YAML. You can sets extra options
with \texttt{natbiboptions} variable in YAML header. Example

\begin{verbatim}
natbiboptions: longnamesfirst,angle,semicolon
\end{verbatim}

There are various more specific bibliography styles available at
\url{https://support.stmdocs.in/wiki/index.php?title=Model-wise_bibliographic_style_files}.
To use one of these, add it in the header using, for example,
\texttt{biblio-style:\ model1-num-names}.

\subsection{Using CSL}\label{using-csl}

If \texttt{cite-method} is set to \texttt{citeproc} in
\texttt{elsevier\_article()}, then pandoc is used for citations instead
of \texttt{natbib}. In this case, the \texttt{csl} option is used to
format the references. By default, this template will provide an
appropriate style, but alternative \texttt{csl} files are available from
\url{https://www.zotero.org/styles?q=elsevier}. These can be downloaded
and stored locally, or the url can be used as in the example header.

\section{Equations}\label{equations}

Here is an equation: \[ 
  f_{X}(x) = \left(\frac{\alpha}{\beta}\right)
  \left(\frac{x}{\beta}\right)^{\alpha-1}
  e^{-\left(\frac{x}{\beta}\right)^{\alpha}}; 
  \alpha,\beta,x > 0 .
\]

Inline equations work as well: \(\sum_{i = 2}^\infty\{\alpha_i^\beta\}\)

\section{Figures and tables}\label{figures-and-tables}

Figure~\ref{fig-meaningless} is generated using an R chunk.

\begin{figure}

\centering{

\includegraphics[width=0.5\linewidth,height=\textheight,keepaspectratio]{arXiv_files/figure-pdf/fig-meaningless-1.pdf}

}

\caption{\label{fig-meaningless}A meaningless scatterplot}

\end{figure}%

\section{Tables coming from R}\label{tables-coming-from-r}

Tables can also be generated using R chunks, as shown in
Table~\ref{tbl-simple} example.

\begin{Shaded}
\begin{Highlighting}[]
\NormalTok{knitr}\SpecialCharTok{::}\FunctionTok{kable}\NormalTok{(}\FunctionTok{head}\NormalTok{(mtcars)[,}\DecValTok{1}\SpecialCharTok{:}\DecValTok{4}\NormalTok{])}
\end{Highlighting}
\end{Shaded}

\begin{longtable}[]{@{}lrrrr@{}}

\caption{\label{tbl-simple}Caption centered above table}

\tabularnewline

\toprule\noalign{}
& mpg & cyl & disp & hp \\
\midrule\noalign{}
\endhead
\bottomrule\noalign{}
\endlastfoot
Mazda RX4 & 21.0 & 6 & 160 & 110 \\
Mazda RX4 Wag & 21.0 & 6 & 160 & 110 \\
Datsun 710 & 22.8 & 4 & 108 & 93 \\
Hornet 4 Drive & 21.4 & 6 & 258 & 110 \\
Hornet Sportabout & 18.7 & 8 & 360 & 175 \\
Valiant & 18.1 & 6 & 225 & 105 \\

\end{longtable}


\renewcommand\refname{References}
  \bibliography{bibliography.bib}



\end{document}
